\chapter{Certified Answers}
\label{ch:certified}

% Written in the v0.507 -> v0.6664 book pass (BOOK_UPDATE_REVIEW.md I27).

Every answer in this book so far has been a \emph{claim}. When \B{Factor} prints
\((x-1)(x+1)(x^2+1)\), Mathilda is asserting that the product expands back to
\(x^4-1\) and that no factor splits further. When \B{Minimize} returns \(\{-1,
\{x\to-1\}\}\), it is asserting that no real \(x\) does better. The assertions are
made by a few hundred thousand lines of C and Mathilda-language rules. That code is
tested, and the earlier chapters show it doing a great deal correctly, but it is
still a program, and programs have bugs. This chapter is about the other way of
trusting an answer: having it checked by a second program that is small, simple,
and independent of the first.

Since version 0.6536 Mathilda can hand its answers to the
\emph{Lean~4}\index{Lean 4} theorem prover. \mcode{Certify}\bidx{Certify} computes a
result, writes a self-contained Lean file that states the result as a theorem and
proves it from a \emph{certificate}, and runs the Lean kernel on that file. The
answer comes back as a \mcode{ProofObject}\bidx{ProofObject} that records exactly what
was proved, how, and what was \emph{not} checked. Around \mcode{Certify} sit
\mcode{LeanForm}\bidx{LeanForm} (Mathilda expressions in Lean syntax),
\mcode{Refute}\bidx{Refute} (a counterexample, with a proof that it is one),
\mcode{Conjecture}\bidx{Conjecture} (guess a closed form, then certify it),
\mcode{LeanVerify}\bidx{LeanVerify} (check a Lean source string), and a small Lean
package that lets a Lean user call Mathilda from inside a proof.%
\calloutnote{Pitfall}{The machine that built this edition of the book has no Lean
toolchain installed. That is deliberate, and it is worth seeing: everything in this
chapter that does not need Lean really runs, and every place where Lean would be
consulted says so honestly, with \mcode{Missing["NotAvailable"]} rather than a
made-up \mcode{True}. Where the chapter shows Lean code that Mathilda did not print,
it is in a tinted box labelled \emph{not build-verified}.}

The chapter starts with why certificates matter and how they relate to the evidence
Mathilda already produces for itself. It then takes each function in turn, describes
the architecture and the trust boundary, shows the Lean side, and ends with an honest
list of limits.

% ---------------------------------------------------------------------------
\section{Why certificates}
\label{sec:cert-why}

There is a basic asymmetry in mathematics between \emph{finding} and
\emph{checking}. Finding the factors of \(2^{64}+1\) was a celebrated hand computation by Landry in 1880;
checking that \(274177 \times 67280421310721 = 18446744073709551617\) takes a
pocket calculator a moment. Finding an antiderivative needs the Risch machinery of
Chapter~\ref{ch:integration}; checking one needs only differentiation. Finding the
minimum of a polynomial needs a search; checking that the polynomial minus the
claimed minimum is a sum of squares needs only multiplication.

A \emph{certificate}\index{certificate} is the extra data that makes checking cheap:
the factors themselves, a list of squares, a set of cofactors, a primality witness.
The hard, clever, bug-prone search produces the certificate. A short, dull, reliable
procedure checks it. If the check passes, the answer is right \emph{regardless of how
the search was done}, and a bug in the search can cost you an answer but never give
you a wrong one.%
\calloutnote{Theory}{This is the \emph{de Bruijn criterion}\index{de Bruijn
criterion}, named after the designer of the Automath proof checker: a proof system is
trustworthy when its proofs can be checked by a small, independently auditable
program, however large and clever the system that \emph{produced} them. Lean~4
follows it. Its \emph{kernel}\index{Lean 4!kernel} type-checks every definition and
theorem against the rules of dependent type theory, and everything else (the
elaborator, the tactics such as \texttt{grind} and \texttt{decide}, the libraries) only
proposes proof terms that the kernel must accept.}

Mathilda adopts the same division of labour. In the language of the design
note\index{Lean 4!trust model}, Mathilda is an \emph{untrusted oracle} and the Lean
kernel is the \emph{trusted checker}. Mathilda's algorithms (Risch, Gröbner bases,
semidefinite programming, factoring) are never trusted; they only have to be right
often enough to be useful. A bug in Mathilda can make \mcode{Certify} fail to find a
certificate, or produce one that Lean rejects. It cannot produce a wrong
\mcode{"Verified" -> True}, except through steps that the \mcode{ProofObject} lists
explicitly as trusted, or through a bug in Lean itself.

% ---------------------------------------------------------------------------
\section{The certificates Mathilda already makes}
\label{sec:cert-internal}

Certificates are not new to Mathilda. Several subsystems described elsewhere in this
book already produce evidence for their answers and check it exactly before they
answer. What is new is that the evidence can now leave the system.

\begin{itemize}
\item \B{PossibleZeroQ} has a \emph{structural non-zero certificate} that settles
  many \texttt{False} answers with no numeric work, and its \texttt{True} answers are
  evidence, not proofs (\S\ref{ssec:alg-possiblezeroq}). When certainty is needed,
  \B{Simplify} can sometimes reduce a difference to the literal \(0\) by an exact
  argument, which \emph{is} a proof.
\item \B{Reduce} tries a sum-of-squares (Putinar) certificate before running the
  cylindrical decomposition, and re-verifies the rounded certificate in exact rational
  arithmetic (\S\ref{ssec:alg-reduce}).
\item The exact optimizers \B{Minimize} and \B{Maximize} finish with a certificate:
  a \B{Reduce} decision of \(\forall x\,(\text{constraints}\Rightarrow f(x)\ge v)\)
  (\S\ref{ssec:opt-exact-method}).
\item \B{Refine}'s sign reasoning is backed by a \B{Reduce} entailment
  (\S\ref{ssec:alg-refine}).
\item The Risch decision procedure can \emph{prove} that an integral is not
  elementary (Chapter~\ref{ch:integration}); the proof is the failure of a complete
  algorithm.
\end{itemize}

The sum-of-squares engine behind \B{Reduce} is exposed directly, which makes the idea
concrete.

\begin{codepairs}
\pair{certified-answers/internal/1}{\mcode{Reduce\textasciigrave SOSCertificate}
proves \(q\ge0\) by returning blocks \(\{m, b, Q\}\): a multiplier \(m\), a basis of
polynomials \(b\), and an exact rational positive semidefinite Gram matrix \(Q\).}
\pair{certified-answers/internal/2}{Anyone can check it: \(q=\sum m\,b^{\mathsf T}Qb\)
must hold identically. The difference expands to \(0\).}
\pair{certified-answers/internal/3}{\(x^2-1\) is negative at \(x=0\), so there is no
certificate, and the engine says so.}
\pair{certified-answers/internal/4}{The zero test's \texttt{True} here is a strong
verdict but, as \S\ref{ssec:alg-possiblezeroq} explains, not a proof.}
\pair{certified-answers/internal/5}{The exact optimizer's answer rests on an internal
\B{Reduce} certificate.}
\end{codepairs}

All of these share one weakness: the certificate is checked by \emph{the same system
that produced it}. Exact rational arithmetic is a strong check, but if
\B{Expand} had a bug, the check and the computation would share it. \mcode{Certify}
closes that gap by sending the certificate to a checker with nothing in common with
Mathilda. Several of its generators reuse the internal machinery directly: the
nonnegativity kind takes its squares from \mcode{Reduce\textasciigrave
SOSCertificate}, and the optimality kind takes the optimum from \B{Minimize} or
\B{Maximize} and then builds its own Lagrange-multiplier certificate. Others have no
external form yet. \B{PossibleZeroQ} answers are never certified (they are
heuristic), and Risch non-elementarity proofs are not exported: a proof that
\emph{no} elementary antiderivative exists is a statement about differential fields
that core Lean cannot yet state (\S\ref{sec:cert-limits}).

% ---------------------------------------------------------------------------
\section{Speaking Lean: \texttt{LeanForm}}
\label{sec:cert-leanform}

Before Mathilda can ask Lean anything it has to say it in Lean. \mcode{LeanForm} is
a display wrapper, like \mcode{FortranForm}: it prints an expression as Lean~4
syntax, with Lean's precedences (\(\leftrightarrow\) 20, \(\to\) 25, \(\lor\) 30,
\(\land\) 35, relations 50, \texttt{+ -} 65, \texttt{* /} 70, \texttt{\^{}} 75),
Lean's explicit \texttt{*}, and Mathlib's names for the elementary functions.

\begin{codepairs}
\pair{certified-answers/leanform/1}{Multiplication is explicit and \texttt{==} becomes
Lean's \texttt{=}.}
\pair{certified-answers/leanform/2}{Negative powers are division, a rational
coefficient is parenthesised, and negative terms print as subtraction.}
\pair{certified-answers/leanform/3}{\B{Function} is \texttt{fun}, application is
juxtaposition with compound arguments in parentheses, and \B{Log} and \B{Pi} are
Mathlib's \texttt{Real.log} and \texttt{Real.pi}.}
\pair{certified-answers/leanform/4}{An inequality chain becomes a conjunction, and
\mcode{Not} is \(\neg\).}
\end{codepairs}

A display form may be approximate. A theorem statement may not: if the Lean text
means something slightly different from the Mathilda expression, Lean proves the wrong
theorem. The certificate generators therefore use a \emph{strict
translator}\index{LeanForm@\mcode{LeanForm}!strict translator},
\mcode{Lean\textasciigrave Internal\textasciigrave ToLean[expr, type, coreOnly,
heads, closed]}, which refuses every node whose Lean reading would differ. Each
refusal issues \mcode{LeanForm::untrans} with the reason and returns
\mcode{\$Failed}.

\begin{codepairs}
\pair{certified-answers/leanform/5}{Over \texttt{Int}, Lean's \texttt{/} is floor
division, so \(x/2\) is refused (``division over Int'').}
\pair{certified-answers/leanform/6}{An inexact number has no faithful Lean literal:
``inexact number; Rationalize it first''.}
\pair{certified-answers/leanform/7}{With \texttt{coreOnly} set, \B{Log} needs Mathlib,
the Tier~1 library, and is refused.}
\pair{certified-answers/leanform/8}{In closed mode a variable not bound by a quantifier
is refused: a free \texttt{id} or \texttt{default} would silently mean Lean's own
global.}
\end{codepairs}

\calloutnote{Pitfall}{The most dangerous difference is division by zero. In Lean,
\(x/0=0\) by definition, for every field. A statement that divides by a quantity that
can vanish is therefore not the statement a mathematician reads. This is why
\mcode{Certify} clears denominators before stating a rational identity, why it lists
the denominators in \mcode{"Hypotheses"}, and why a \B{Solve} certificate states every
denominator nonzero at every solution: otherwise Lean's \(x/0=0\) would accept the
spurious root \(x=1\) of \((x-1)/(x^2-1) = (x-1)/(2x-2)\).}

% ---------------------------------------------------------------------------
\section{Certify}
\label{sec:cert-certify}

\mcode{Certify[comp]} evaluates a computation and certifies its result;
\mcode{Certify[a == b]} certifies a claim; \mcode{Certify[\{c1, c2, \ldots\}]}
certifies a batch in a single Lean run. The result is a \mcode{ProofObject}, an
association with a fixed set of properties.

\subsection{A first certificate}

\begin{codepairs}
\pair{certified-answers/first/1}{No Lean executable was found on this machine.
\mcode{\$LeanExecutable} is \texttt{Automatic}: it tries \texttt{MATHILDA\_LEAN},
the newest \texttt{\textasciitilde/.elan} toolchain, then \texttt{lean} on the
\texttt{PATH}.}
\pair{certified-answers/first/2}{Certify a factorization. The Lean file is written and
Lean is called; since it cannot be run, the message \mcode{LeanVerify::nolean} is
issued (``no Lean 4 executable found'').}
\pair{certified-answers/first/3}{So the verdict is \mcode{Missing["NotAvailable"]}. A
missing checker is not a refutation, so it is never \texttt{False}.}
\pair{certified-answers/first/4}{The certificate kind.}
\pair{certified-answers/first/5}{\mcode{"Result"} is \B{Factor}'s own answer.}
\pair{certified-answers/first/6}{\mcode{"Theorem"} is the statement handed to Lean, as
a held Mathilda expression: for every \(x\), \(x^4-1\) equals the product.}
\pair{certified-answers/first/7}{What Lean does \emph{not} check. The theorem proves
that the product is right, not that the factors are irreducible.}
\pair{certified-answers/first/8}{The full property list. \mcode{"Axioms"},
\mcode{"Logs"} and \mcode{"LeanVersion"} report what the Lean run said.}
\end{codepairs}

The property \mcode{"TrustedSteps"}\index{certificate!trusted steps} is the heart of
the design. It is a precise statement of the gap between what was proved and what was
claimed. For a polynomial identity it is empty: the theorem \emph{is} the claim. For a
factorization it says that irreducibility is not proved. A reader who needs the
irreducibility as well knows exactly what is missing.

\subsection{What Lean is given}

\mcode{"Verify" -> False} builds the whole Lean file without running anything. It is
the way to look at a certificate on a machine without Lean, and the way to take the
file to a machine that has it.

\begin{codepairs}
\pair{certified-answers/source/1}{Generate, but do not check.}
\pair{certified-answers/source/2}{\mcode{Missing["NotChecked"]}: nothing was run.}
\pair{certified-answers/source/3}{The statement in Lean. All coefficients are
integers and there is no division, so the theorem is over \texttt{Int}.}
\pair{certified-answers/source/4}{The proof script, split into its lines: introduce
the variable, then let \texttt{grind} decide the ring identity.}
\pair{certified-answers/source/5}{The end of the file, after the prelude. The theorem
\texttt{mathilda\_cert\_1}, then a \emph{trusted footer} that restates the statement
and asks Lean for the axioms the proof uses.}
\pair{certified-answers/source/6}{The whole file is about 70 lines: most of it is the
prelude, a fixed block of definitions (\texttt{powMod}, \texttt{prattCheck},
\texttt{matMul}, \texttt{det}, \ldots) inlined into every file.}
\end{codepairs}

\calloutnote{Theory}{Why is \texttt{Int} enough? A polynomial identity with integer
coefficients that holds for every integer holds identically in \(\mathbb{Z}[x]\) (a
nonzero polynomial has finitely many roots), and therefore in every commutative ring,
\(\mathbb{Q}\), \(\mathbb{R}\) and \(\mathbb{C}\) included. Stating it over
\texttt{Int} loses nothing and lets the tactic work in the simplest ring there is.}

The footer is what turns ``Lean printed no error'' into a precise
guarantee\index{Lean 4!what Verified means}. \mcode{"Verified" -> True} means all of
the following:
\begin{enumerate}
\item no message of severity \texttt{error}, and no \texttt{sorry}, lies in the
  theorem's lines or its footer lines, or in the header and prelude;
\item the axiom report printed \emph{on the theorem's own footer line} names only
  \texttt{propext}, \texttt{Classical.choice} and \texttt{Quot.sound}, the three
  standard axioms (\texttt{sorryAx}, and the extra axiom that \texttt{native\_decide}
  introduces, both fail);
\item Lean exited cleanly, or every error it reported belongs to a sibling theorem of
  the same batch;
\item a \emph{kernel replay}\index{Lean 4!kernel replay} accepted it: the file is
  compiled to an \texttt{.olean}, and a second, fresh \texttt{lean} process, running
  none of the checked file's code, re-checks every declaration in a new kernel
  environment and computes the axioms itself.
\end{enumerate}
The replay catches a declaration that was added without a kernel check. Behind it is
a simpler foundation: the generated file contains no user-defined code at all, only
the fixed prelude, the strictly translated statement, and a proof script
(\S\ref{sec:cert-leanverify}).

\subsection{The claim as written}
\label{ssec:cert-as-written}

A certificate is only as good as the faithfulness of its statement. The subtle danger
is that a CAS evaluates things. If \mcode{Certify} evaluated the claim
\mcode{Simplify[(x\^{}2 - 1)/(x - 1)] == x + 1} before translating it, the left side
would become \(1+x\) and Lean would prove the tautology \(1+x=1+x\), which certifies
nothing about \B{Simplify}. So \mcode{Certify} and \mcode{ProofObject} hold their
arguments, and the claim is walked \emph{held}\index{certificate!claim as written}.
Arithmetic, relations, logic, quantifiers and the value-preserving heads \B{Expand},
\B{Together}, \B{Cancel}, \B{Apart}, \B{Factor} and \B{Collect} are read through; any
other computation is refused.

\begin{codepairs}
\pair{certified-answers/aswritten/1}{Refused, with \mcode{Certify::eval}: the claim
applies \B{Simplify}, which Mathilda would have to evaluate before Lean sees it.
Certify the computation \mcode{Certify[Simplify[\ldots]]} instead.}
\pair{certified-answers/aswritten/2}{The same claim with arithmetic only.}
\pair{certified-answers/aswritten/3}{Denominators are cleared \emph{as written},
without cancellation: Lean proves \(x^2-1=(x-1)(1+x)\), not \(1+x=1+x\).}
\pair{certified-answers/aswritten/4}{The side condition, \(x\ne1\).}
\pair{certified-answers/aswritten/5}{And the one trusted step: the field algebra from
the cross-multiplied identity back to the rational one, where the denominators are
nonzero.}
\pair{certified-answers/aswritten/6}{Give \(x\) a value.}
\pair{certified-answers/aswritten/7}{Refused with \mcode{Certify::value} (``x has a
value (5)''). The claim is certified as written, never with values substituted.}
\pair{certified-answers/aswritten/8}{}
\pair{certified-answers/aswritten/9}{A transcendental identity has no Tier~0
generator: \B{Sin} is refused with \mcode{Certify::eval}.}
\pair{certified-answers/aswritten/10}{An identity under a hypothesis.}
\pair{certified-answers/aswritten/11}{It is stated for every field \(K\) of
characteristic~0, so \(\mathbb{Q}\), \(\mathbb{R}\) and \(\mathbb{C}\) at once.}
\pair{certified-answers/aswritten/12}{Hence \mcode{"Domain"} is \mcode{Complexes}.}
\end{codepairs}

\calloutnote{Theory}{Why a generic field when there are hypotheses? Over
\(\mathbb{Q}\) a hypothesis can be unsatisfiable, and an implication with a false
hypothesis is vacuously true. ``For all rational \(x\), \(x^2=-4\Rightarrow x=7\)'' is
a theorem, because no rational \(x\) has \(x^2=-4\); but Mathilda treats \(x\) as
complex, and \(x=2i\) is a counterexample. Stating the claim for every
characteristic-0 field closes the hole: \(\mathbb{C}\) is such a field, so a proof
covers it. When the hypotheses contain an order relation, every variable must be
real, and the statement is made over every ordered field instead
(\mcode{"Domain" -> Reals}).}

Two further details complete the picture. Mathilda puts the held \B{Plus},
\B{Times} and \B{Power} nodes in its canonical arithmetic form before translation
(flattened, ordered, rationals reduced), so Lean checks \((1+x)^2\) where you typed
\((x+1)^2\). That normalisation is part of the trust base, and
\mcode{"LeanTheorem"} shows exactly what was proved. And a \mcode{ProofObject} is
\emph{authenticated}: \mcode{Certify} registers each object it produces by its
\B{Hash}, so a hand-made \mcode{ProofObject[<|"Verified" -> True|>]} reports
\mcode{Missing["Unauthenticated"]} (\S\ref{sec:cert-engines}).

\subsection{A tour of certificate kinds}
\label{ssec:cert-kinds}

Each kind of computation has its own generator, its own certificate, and its own
proof strategy. Table~\ref{tab:cert-kinds} summarises them.

\begin{table}[ht]
\centering\small
\begin{tabular}{@{}p{0.20\linewidth}p{0.27\linewidth}p{0.22\linewidth}p{0.23\linewidth}@{}}
\toprule
Kind & Certificate & Lean proof & Trusted steps \\
\midrule
Identity & the claim itself, denominators cleared & \texttt{grind} & none (polynomial) \\
Factorization & the product & \texttt{grind} & irreducibility \\
Antiderivative & cleared derivative identity, surds as variables & \texttt{grind} & derivative rules, \B{Together} forms \\
Bezout & \(s,t\) and the quotients & \texttt{grind} & none \\
IdealMembership & cofactors \(c_i\) & \texttt{grind} & none \\
Roots & the factorization into linear factors & \texttt{grind} & none if it splits over \(\mathbb{Q}\) \\
Matrix kinds & the matrices & \texttt{decide +kernel} & \B{NullSpace} completeness \\
Nonnegativity & weighted squares \(d_k\ell_k^2 m_k\) & hints, \texttt{grind} & none \\
Optimality & KKT multipliers, squares & hints, \texttt{grind} & none \\
Primality & a Pratt tree & \texttt{decide +kernel} & Lucas's theorem \\
ClosedFormSum & the closed form & induction & integer \(n\ge a\) only \\
\bottomrule
\end{tabular}
\caption{The \mcode{"Kind"}s of \mcode{Certify} at Tier~0 (core Lean, no Mathlib).}
\label{tab:cert-kinds}
\end{table}

\paragraph{Antiderivatives.} \mcode{Certify[Integrate[f, x]]} proves, in core Lean,
that \(f\) equals the derivative of Mathilda's antiderivative \(F\). Core Lean knows
nothing about \(\log\) or \(\arctan\), so the derivative is taken by Mathilda, the
result is cleared of denominators, and Lean checks the polynomial identity that
remains.

\begin{codepairs}
\pair{certified-answers/antideriv/1}{Make \mcode{"Verify" -> False} the session
default, so that each certificate is built but not run.}
\pair{certified-answers/antideriv/2}{A rational integrand whose antiderivative
contains \(\sqrt2\).}
\pair{certified-answers/antideriv/3}{Mathilda's answer.}
\pair{certified-answers/antideriv/4}{The theorem. \(\sqrt2\) has become a fresh
variable \texttt{s1}, and the right side is a multiple of \(s_1^2-2\).}
\pair{certified-answers/antideriv/5}{Four trusted steps: the elementary derivative
rules, Mathilda's \B{D} of the polynomials, the \B{Together} normal forms, and that
the cleared denominators are not identically zero.}
\pair{certified-answers/antideriv/6}{The famous three-arctangent integral, re-checked
by the exact engine (\S\ref{sec:cert-engines}).}
\pair{certified-answers/antideriv/7}{A claimed antiderivative is certified by its
derivative, \(D[F,x]=f\), with the message \mcode{Certify::antideriv}: Lean cannot
know \emph{which} antiderivative \B{Integrate} returns, so the claim
\(\int x\,dx = x^2/2+7\) is true only up to a constant.}
\end{codepairs}

\calloutnote{Theory}{Why is it sound to replace \(\sqrt2\) by a variable \(s\)? The
statement Lean proves is an identity \(P(x,s) = (s^2-2)K(x,s)\) in the polynomial ring
\(\mathbb{Q}[x,s]\), with no assumption about \(s\) at all. Substituting \(s\mapsto\sqrt2\)
is a ring homomorphism into \(\mathbb{R}[x]\), so the identity survives it, and the right
side becomes \(0\). The tempting alternative, ``assume \(s^2=2\) over \(\mathbb{Q}\)'', would
be vacuous, because no rational \(s\) satisfies it, and a theorem with an unsatisfiable
hypothesis proves nothing. Square roots are written over primes
(\(\sqrt{12}=2\sqrt3\), \(\sqrt6=\sqrt2\sqrt3\)) so that the relations between them stay
visible to the identity.}

An antiderivative with a \B{RootSum}, as for \(1/(x^3+x+1)\), is outside the Tier~0
shape (\mcode{Certify::tier}, ``needs Tier 1 (Mathlib)''). With a Lake project with
Mathlib built (\mcode{\$LeanProject}), \mcode{Tier -> 1} proves the analytic statement
\(\forall x\in\mathbb{R},\ d_1\ne0\to\dots\to\texttt{HasDerivAt}\,(\lambda x.\,F)\,f\,x\)
directly, with no trusted steps. Mathilda then also checks that the hypotheses
\(d_i\neq 0\) are satisfiable (it finds a rational witness point), because Lean's
\(\log x = \log|x|\) can make an innocent-looking hypothesis false for every real \(x\),
and the theorem vacuous.

\paragraph{Sums of squares and optimality.} A claim \(\forall x\,(p\ge q)\) is
certified by writing \(p-q=\sum_k d_k\,\ell_k^2\,m_k\), with rational \(d_k>0\),
polynomials \(\ell_k\), and each \(m_k\) either \(1\) or a hypothesis \(g\ge0\). Lean
then needs only the facts \(0\le\ell^2\) and the identity.

\begin{codepairs}
\pair{certified-answers/sos/1}{The arithmetic--geometric mean inequality for two
numbers.}
\pair{certified-answers/sos/2}{The proof: the hint \(0\le(x-y)^2\), the identity
\(x^2-2xy+y^2=(x-y)^2\) by \texttt{grind}, then \texttt{grind} closes the goal.}
\pair{certified-answers/sos/3}{A quartic, checked exactly.}
\pair{certified-answers/sos/4}{Its squares came from \B{Reduce}'s Putinar engine,
the one shown in \S\ref{sec:cert-internal}.}
\pair{certified-answers/sos/5}{A linear program.}
\pair{certified-answers/sos/6}{The optimum, and the exact re-check of its
certificate.}
\pair{certified-answers/sos/7}{One Lagrange multiplier per inequality. The two
inactive constraints get \(0\).}
\pair{certified-answers/sos/8}{The leftover \(\sigma\) is zero, so the problem is an LP
and the multipliers are its dual certificate.}
\end{codepairs}

\calloutnote{Theory}{An optimality certificate has two halves. Attainment is a
calculation: the point \(x^*\) is feasible and \(f(x^*)=m\). The bound needs
multipliers \(\lambda_j\ge0\) such that \(\sigma = f-m-\sum_j\lambda_j g_j\) is a sum
of squares; then on the feasible set every \(g_j\ge0\), so \(f-m\ge\sigma\ge0\). For a
linear program \(\sigma\equiv0\) and this is weak LP duality. The Gram matrices for
\(\sigma\) are found numerically, but the certificate Lean sees is exact: the matrix is
factored by an exact rational \(LDL^{\mathsf T}\), \(Q = L\,\mathrm{diag}(d)\,L^{\mathsf T}\),
and \(\ell = L^{\mathsf T} b\). Both kinds are stated over \emph{every ordered field},
which includes \(\mathbb{R}\), so nothing about the reals is left unchecked.}

\paragraph{Primes.} A prime is certified by a \emph{Pratt
certificate}\index{Pratt certificate}\index{primality!certificate}: a witness \(a\)
with \(a^{p-1}\equiv1\) and \(a^{(p-1)/q}\not\equiv1 \pmod p\) for every prime \(q\)
dividing \(p-1\), together with Pratt certificates for those \(q\), recursively. The
prelude's \texttt{Mathilda.prattCheck} checks such a tree, and Lean's kernel evaluates
it with \texttt{decide +kernel}.

\begin{codepairs}
\pair{certified-answers/primes/1}{Landry's factorization of the Fermat number
\(F_6\).}
\pair{certified-answers/primes/2}{The theorem: the product is right, and each factor
has a Pratt certificate.}
\pair{certified-answers/primes/3}{The exact engine replays the certificates with
\B{PowerMod}.}
\pair{certified-answers/primes/4}{A composite is certified by a divisor, against the
prelude's definition of a prime. This one needs no trusted step.}
\pair{certified-answers/primes/5}{The flat certificate for \(274177\), in dependency
order: \((p, a, \text{factors of } p-1)\), with \(2\) as the base.}
\end{codepairs}

\calloutnote{Theory}{Lucas's theorem says that if \(a^{p-1}\equiv 1\) and
\(a^{(p-1)/q}\not\equiv 1\pmod p\) for every prime \(q\mid p-1\), then \(a\) has order
\(p-1\) in \((\mathbb{Z}/p)^\times\), so that group has \(p-1\) elements and \(p\) is
prime. Pratt observed (1975) that recursing on the \(q\) gives a certificate of
polynomial size, which shows that primality is in NP. Lean checks the certificate;
it does not prove Lucas's theorem, which is therefore listed as a trusted step.
The kernel's \texttt{powMod} unfolds one step per bit of the exponent, so even a
prime of more than 260 bits, such as \(51\cdot2^{257}+1\), checks in well under a
second.}

\paragraph{Sums, roots, matrices and ideals.} The remaining kinds follow the same
pattern.

\begin{codepairs}
\pair{certified-answers/kinds/1}{Build, do not run.}
\pair{certified-answers/kinds/2}{\(\sum_{k=0}^n k\,2^k\): Lean defines the partial sum
\texttt{mathilda\_cert\_1\_S} by recursion and proves the closed form by induction on
\(m\), with \(n=m\).}
\pair{certified-answers/kinds/3}{A cubic that splits over \(\mathbb{Q}\): the
certificate proves \emph{completeness}, an ``if and only if'', not just that the three
roots work.}
\pair{certified-answers/kinds/4}{A matrix inverse. The shapes are stated too, because
the prelude's list functions would silently truncate ragged input.}
\pair{certified-answers/kinds/5}{An implication between polynomial equations.}
\pair{certified-answers/kinds/6}{The cofactors: \(2y^2-1 = 1\cdot(x^2+y^2-1)
+(-x-y)(x-y)\).}
\pair{certified-answers/kinds/7}{Stated over every field of characteristic~0.}
\end{codepairs}

\calloutnote{Under the hood}{The ideal-membership cofactors come from
\B{PolynomialReduce} when its remainder is \(0\), and otherwise from a degree-bounded
ansatz (degrees \(0\) to \(4\)) whose coefficient equations are solved exactly by
\B{LinearSolve}. In Lean the proof is: \texttt{have} \(k_i: l_i - r_i = 0\) from each
hypothesis, \texttt{have} the cofactor identity by \texttt{grind}, rewrite with the
\(k_i\), and close with \texttt{grind}. The \B{Solve} kind certifies completeness only
for a single univariate equation that splits into linear factors over \(\mathbb{Q}\);
for systems, constraints and irreducible factors it certifies soundness (every
solution satisfies every constraint) and lists completeness as a trusted step.}

% ---------------------------------------------------------------------------
\section{Engines and consensus}
\label{sec:cert-engines}

Lean is one \emph{engine}\index{certificate!engine}. The option \mcode{"Engine"}
chooses among three, and each declares how much it may be trusted.

\begin{description}
\item[\texttt{"Lean"}] (the default) is a kernel (\texttt{"Trust" -> "Kernel"}). It is
  the only engine whose verdict can be \texttt{True}.
\item[\texttt{"Exact"}] re-checks the certificate with Mathilda's own exact rational
  arithmetic and polynomial \B{Expand}, with no Lean. It recomputes the obligations from
  the certified statement, never from the generator's own pre-check. Because it is the
  same system that produced the certificate, its verdicts are
  \mcode{Missing["Exact", True]} or \mcode{Missing["Exact", False]}, never \texttt{True}.
\item[\texttt{"Mathilda"}] is a heuristic check through \B{PossibleZeroQ} and
  \B{Simplify}, reported as \mcode{Missing["Heuristic", \ldots]}.
\end{description}

The labels are enforced, not merely declared: a \texttt{True} from any backend whose
trust is not \texttt{"Kernel"} is rewritten to \mcode{Missing[name, True]} before it
reaches a \mcode{ProofObject}.

\begin{codepairs}
\pair{certified-answers/engines/1}{The exact engine, on a machine without Lean.}
\pair{certified-answers/engines/2}{Every check passed.}
\pair{certified-answers/engines/3}{The checks themselves: the cleared identity expands
to \(0\), and the cleared denominator is not identically zero.}
\pair{certified-answers/engines/4}{A false identity (\(2x\) should be \(3x\)). The
exact check fails, with \mcode{Certify::exactfail} (quieted here).}
\pair{certified-answers/engines/5}{The heuristic engine's answer is labelled as such.}
\pair{certified-answers/engines/6}{A list of engines runs them all.}
\pair{certified-answers/engines/7}{Each engine's verdict.}
\pair{certified-answers/engines/8}{The combined verdict. A kernel was listed but could
not run, so the verdict is the kernel's: not available.}
\pair{certified-answers/engines/9}{With no kernel listed, agreement is reported as
\mcode{Missing["Consensus", True]}: two checks agree, but neither is a proof.}
\pair{certified-answers/engines/10}{A forged object.}
\pair{certified-answers/engines/11}{It was not produced by \mcode{Certify} in this
session, so its claimed verdict is not reported.}
\end{codepairs}

The consensus rule\index{certificate!consensus} for a list of engines is
conservative. Any \texttt{False} or \mcode{Missing[\_, False]} makes the result
\texttt{False}, and if a kernel said \texttt{True} at the same time,
\mcode{Certify::disagree} reports that some checker or translation is wrong. Otherwise
the result is \texttt{True} only when every kernel listed said \texttt{True}; with no
kernel listed it is at best \mcode{Missing["Consensus", True]}.

A second engine earns its place when it disagrees. When the Lean integration landed,
the following near miss was its showcase: \B{PossibleZeroQ} called the difference
\((1+x)^{20} - \operatorname{Expand}[(1+x)^{20}] + 10^{-25}x\) zero, the heuristic
engine agreed, and Lean refuted the identity. The zero test has since been fixed, and
today every engine on this machine gets it right.

\begin{codepairs}
\pair{certified-answers/catches/1}{A polynomial that is exactly \(x/10^{25}\).}
\pair{certified-answers/catches/2}{The zero test now sees that it is not zero.}
\pair{certified-answers/catches/3}{The claim that it \emph{is} an identity, checked by
two non-kernel engines. Both reject it. \B{Expand} is one of the value-preserving
heads that the claim walker reads through, so the claim is checked as written.}
\end{codepairs}

\calloutnote{Pitfall}{\mcode{Missing["Exact", True]} is a strong answer. It means that
an exact rational computation, independent of the code path that produced the
answer, came out right. But it is a check by the same system, sharing its arithmetic
and its \B{Expand}, and the book reports it as what it is. Only a run of the Lean
kernel turns it into \texttt{True}.}

% ---------------------------------------------------------------------------
\section{Refute}
\label{sec:cert-refute}

A false claim deserves a certificate too. \mcode{Refute[claim]} searches for a
counterexample and then proves the \emph{negation} of the statement \mcode{Certify}
would have tried to prove. The witness is the certificate, and the Lean proof is short:
instantiate the universal statement at the witness and let the kernel compute that the
result is false.

\begin{codepairs}
\pair{certified-answers/refute/1}{Refute a wrong expansion of \((x+1)^3\) (generate
only).}
\pair{certified-answers/refute/2}{The witness.}
\pair{certified-answers/refute/3}{The theorem is the negation of the statement
\mcode{Certify} would prove.}
\pair{certified-answers/refute/4}{Its proof: apply the hypothesis \(h\) at \(x=-1\)
and evaluate.}
\pair{certified-answers/refute/5}{An inequality that fails at the origin.}
\pair{certified-answers/refute/6}{A pole is not a counterexample: \(x=0\) falsifies the
cleared statement but the claim is undefined there, so the witness is \(x=-1\).}
\pair{certified-answers/refute/7}{A hypothesis with no rational solution.}
\pair{certified-answers/refute/8}{The witness \(x=-2i\) is algebraic. Core Lean has no
kernel proof for it, so the refutation is exact-only.}
\pair{certified-answers/refute/9}{The witness must satisfy the assumptions too: with
\(x>0\), the counterexample is \(x=1/2\).}
\pair{certified-answers/refute/10}{No witness within the search bounds:
\mcode{Refute::none}, ``this is not a proof of the claim''. (The polynomial is
Motzkin's, which is nonnegative but not a sum of squares.)}
\end{codepairs}

\calloutnote{Theory}{A witness refutes a statement only if it lies in the statement's
quantifier domain\index{Refute@\mcode{Refute}!faithfulness}. An integer witness
refutes a statement over \texttt{Int}; a rational witness refutes ``for every field of
characteristic~0'', because \(\mathbb{Q}\) is such a field, and also the user's claim
over \(\mathbb{C}\), because \(\mathbb{Q}\subset\mathbb{C}\). The search is
deterministic: an exact grid by increasing height (integers \(-5\ldots5\), or small
fractions), then \B{FindInstance} over the integers or rationals, and finally, exact
only, \B{FindInstance} over \(\mathbb{C}\) or \(\mathbb{R}\). Every candidate is checked by
exact arithmetic before a refutation is built; \B{FindInstance}'s domain argument is
not trusted.}

\mcode{Certify} uses the same machinery on its own. When Lean \emph{rejects} a
certificate, \mcode{Certify} searches for a counterexample for up to two seconds (the
option \mcode{"Counterexample"}); if it finds one, the object gains
\mcode{"Counterexample"}, \mcode{"Refuted"} and \mcode{"RefutationTheorem"}, and is
displayed as \emph{Refuted} rather than merely \emph{not verified}.

% ---------------------------------------------------------------------------
\section{Conjecture}
\label{sec:cert-conjecture}

Experimental mathematics works in three steps: compute some values, guess a pattern,
prove it. \mcode{Conjecture[expr, n]} does all three. It evaluates \texttt{expr}
exactly at \(n = s, s+1, \ldots\) for twelve fit points and forty test points, guesses a
closed form from the fit points, checks it on the test points, and hands the claim
\(\texttt{expr} = \text{guess}\) to \mcode{Certify} with the engines Lean and Exact.

\begin{codepairs}
\pair{certified-answers/conjecture/1}{The sum of fourth powers.}
\pair{certified-answers/conjecture/2}{The guess: Faulhaber's formula, factored.}
\pair{certified-answers/conjecture/3}{It was fitted on \(n=1\ldots12\) and tested on
\(n=1\ldots52\).}
\pair{certified-answers/conjecture/4}{\texttt{"Checked"}: the certificate passed the
exact re-check, but no kernel ran. With Lean installed it would be
\texttt{"Proved"}.}
\pair{certified-answers/conjecture/5}{The two engines' verdicts.}
\pair{certified-answers/conjecture/6}{A polynomial times a geometric term.}
\pair{certified-answers/conjecture/7}{Found as the solution of a linear recurrence.}
\pair{certified-answers/conjecture/8}{A Fibonacci identity.}
\pair{certified-answers/conjecture/9}{\(\sum F_k^2 = F_nF_{n+1}\), certified with a
Lean definition of the Fibonacci numbers.}
\pair{certified-answers/conjecture/10}{The Vandermonde-type sum
\(\sum_k\binom nk^2\).}
\pair{certified-answers/conjecture/11}{The guess \(\binom{2n}{n}\) is right, but there
is no certificate kind for a binomial sum yet, so it is only
\texttt{"Tested"}, on the points in \mcode{"TestedOn"}.}
\end{codepairs}

The four statuses\index{Conjecture@\mcode{Conjecture}!status} are graded by
evidence: \texttt{"Proved"} (a kernel accepted), \texttt{"Checked"} (no kernel
verdict, but the exact re-check passed), \texttt{"Refuted"} (the certification found a
counterexample), and \texttt{"Tested"} (anything else: the guess agrees with the data,
and nothing more is claimed).%
\calloutnote{Under the hood}{The guessers are tried in order. \texttt{"Polynomial"}
interpolates with \B{InterpolatingPolynomial} (degree at most 10).
\texttt{"Geometric"} finds the smallest linear recurrence with constant coefficients
(order at most 6) from the exact \B{NullSpace} of a Hankel matrix of the values; the
rational roots \(r\) of its characteristic polynomial, with multiplicities \(\mu\),
give the ansatz \(\sum_r p_r(n)\,r^n\) with \(\deg p_r<\mu\). \texttt{"Fibonacci"},
\texttt{"Rational"} (\(p(n)/q(n)\) from an exact null space) and
\texttt{"Hypergeometric"} (central binomials, factorials, or a term ratio with rational
roots) follow. A guess is accepted only if it reproduces every fit value
\emph{exactly}; testing on further points is a guard against overfitting, but only the
certificate turns it into a theorem. The closed-form-sum certificate itself is an
induction: Lean checks the base case and the step \(C(n+1)-C(n)=t(n+1)\).}

% ---------------------------------------------------------------------------
\section{LeanVerify}
\label{sec:cert-leanverify}

\mcode{LeanVerify[src]} checks an arbitrary Lean~4 source string and returns an
association: \mcode{"Success"}, Lean's \mcode{"ExitCode"} and \mcode{"Messages"},
the axiom reports, the timing, the command run, and the result of the kernel replay.
It is the tool for a certificate that has been saved and carried to another session,
where its \mcode{ProofObject} would no longer be authenticated: check
\mcode{obj["LeanSource"]} again.

What makes \mcode{LeanVerify} interesting is what it refuses to run. A Lean file can
contain code that runs \emph{while it is being elaborated}: a macro, a custom tactic, a
\texttt{\#eval}. Such code could write the very \texttt{.olean} that the replay reads,
and forge any result. So before Lean starts, the source is checked against a minimal
surface\index{LeanVerify@\mcode{LeanVerify}!source check}, in three layers:
\begin{description}
\item[Lexical] no strings, no comments, no character literals, no \texttt{\#}, no
  \texttt{\$}, \texttt{\%}, backquotes or backslashes, no \texttt{@[}, and only an
  allowlist of Unicode. With these gone, Lean's lexer skips nothing, so the checker
  does not have to reproduce it (\mcode{LeanVerify::lex}).
\item[Words] the command keywords of Lean's parser, derived from its parser tables and
  tested against them, and the escape words such as \texttt{native\_decide},
  \texttt{run\_tac}, \texttt{by\_elab}, \texttt{sorry} and \texttt{unsafe}
  (\mcode{LeanVerify::unsafe}).
\item[Grammar] no imports (or only \texttt{import Mathlib} with a project); commands
  only \texttt{theorem}, \texttt{example}, \texttt{section}, \texttt{end},
  \texttt{variable}, \texttt{universe} and a few numeric or Boolean
  \texttt{set\_option}s (\mcode{LeanVerify::import}, \mcode{LeanVerify::cmd}).
\end{description}
These checks are Mathilda's, so they run without Lean.

\begin{codepairs}
\pair{certified-answers/leanverify/1}{A correct Lean theorem, which passes the source
check.}
\pair{certified-answers/leanverify/2}{Lean itself cannot run here. The shell reports
``command not found'' (exit code 127, with \mcode{LeanVerify::nolean}), the replay is
\texttt{"not run"}, and \mcode{"Success"} is \texttt{False}: nothing was checked.}
\pair{certified-answers/leanverify/3}{\texttt{native\_decide} trusts compiled code
instead of the kernel: refused before Lean would run.}
\pair{certified-answers/leanverify/4}{Imports are refused: \texttt{import Lean} would
bring meta-programming into scope.}
\pair{certified-answers/leanverify/5}{A user definition is refused.}
\pair{certified-answers/leanverify/6}{So is a comment, which Lean's lexer would skip.}
\pair{certified-answers/leanverify/7}{A refused source is never handed to the shell:
its exit code is \mcode{Missing["NotRun"]}, not 127.}
\end{codepairs}

On a machine with Lean, the first input returns \mcode{"Success" -> True}. The
command it runs, for core Lean, has this shape (illustrative, with the temporary
directory abbreviated; it is not output from this build):

\begin{shell}
lean --json -R DIR -o DIR/MathildaCert.olean DIR/MathildaCert.lean
lean --run src/internal/lean/Replay.lean DIR/MathildaCert.olean
\end{shell}

\calloutnote{Under the hood}{Strict C99 has exactly one way to start a process,
\texttt{system()}, and that is what \texttt{src/lean\_run.c} uses. No user text ever
reaches the shell: the command line holds only executable paths, flags written by
Mathilda, and temporary-file paths, each single-quoted for the POSIX shell. The Lean
source, the claim and any tactic string go into the \texttt{.lean} file. Each run gets a
private directory from \texttt{mkdtemp} (mode \texttt{0700}); the source is created in it
with \texttt{O\_CREAT|O\_EXCL|O\_NOFOLLOW}, so no other user can plant a file or a
symbolic link there, and the directory is removed afterwards. A Lean run is bounded by
the deterministic \texttt{maxHeartbeats} of each theorem (the option
\mcode{TimeConstraint}, 60 by default), not by a wall clock.}

% ---------------------------------------------------------------------------
\section{The architecture and the trust boundary}
\label{sec:cert-architecture}

\mcode{Certify} does not know Lean. Each generator emits a prover-neutral
intermediate form\index{certificate!prover-neutral IR}: the kind, the theorem, the
certificate data, the trusted steps, and a typed \emph{obligation}
(\texttt{"Identity"}, \texttt{"WeightedSOS"}, \texttt{"KKT"}, \texttt{"Primality"},
\texttt{"Induction"}, \ldots) together with a \emph{domain form} that says over what the
statement quantifies (\texttt{"Int"}, \texttt{"Rat"}, \texttt{"CharZeroField"},
\texttt{"OrderedField"}, \ldots). Registered backends\index{certificate!backend} then
encode and run the obligations, in this order:

\begin{ccode}[Pipeline sketch, after docs/design/lean\_integration.md]
Certify[c, opts]
  per item:   held claim walker -> generator -> neutral IR (Kind, Theorem,
              Certificate, TrustedSteps, Obligation, DomainForm)
              each engine's "Encode" hook (Lean text, exact checks)
  per engine: "Run" on the whole batch (Lean: cache, one file, lean --json,
              axiom reports, kernel replay)
              guard: a True from a non-kernel engine -> Missing[engine, True]
  per item:   one engine -> its ProofObject; a list -> consensus
              kernel False -> counterexample search -> kernel-checked "not S"
\end{ccode}

The package is written in Mathilda's own language, in \texttt{src/internal/lean/}:
\texttt{Lean.m} is the prover-neutral core (claim walker, one generator per kind, the
backend registry, \mcode{ProofObject}); \texttt{LeanBackend.m} the Lean backend and
\mcode{LeanVerify}; \texttt{ExactBackend.m} the exact re-check; \texttt{Engines.m} the
heuristic engine and consensus; \texttt{Refute.m} and \texttt{Conjecture.m} the last two
functions. Only the translator \mcode{LeanForm} (\texttt{src/print\_lean.c}), the
process runner (\texttt{src/lean\_run.c}) and the source checker are C. This follows
the book's recurring split between C for performance and rules for mathematics: the
generators only orchestrate builtins and assemble text. A new backend (an SMT solver,
say) is one more file that registers \texttt{"Name"}, \texttt{"Trust"},
\texttt{"Encode"}, \texttt{"Run"} and \texttt{"Present"} hooks, and it must prove only
what its domain form says: an SMT solver over the reals proves ``for all reals'', not
``for every ordered field''.

The trust boundary\index{Lean 4!trust base} can now be drawn precisely.

\paragraph{Trusted.} The Lean kernel; the Lean toolchain as installed (elaborator,
compiler, \texttt{.olean} writer and reader); the imported modules as built (core
\texttt{Init}, and Mathlib at Tier~1); the source check, whose two parts are either
trivially inspectable (the minimal lexical surface) or tested against Lean's own
parser tables (the keyword list); the assumption that Lean's built-in tactics do not
run user-supplied code through ordinary syntax; Mathilda's canonical arithmetic form
of the claim; and, per certificate, the steps listed in \mcode{"TrustedSteps"}.

\paragraph{Checked.} Everything else: every algorithm that produced the answer and
the certificate (\B{Factor}, \B{Integrate}, \B{Reduce}'s semidefinite programming,
\B{PrimitiveRoot}, \B{Minimize}), the proof script, and the arithmetic of the
certificate itself. Mathilda cannot make Lean accept a false theorem. It can make the
theorem differ from the claim only through the listed trusted steps and the
arithmetic normalisation, and \mcode{"Theorem"} and \mcode{"LeanTheorem"} show exactly
what was proved.%
\calloutnote{Performance}{A Tier~0 run takes about half a second of Lean time plus
about a second for the kernel replay; the design note's whole-call figures (on a
loaded machine, not idle) are about 1.8~s for an accepted \mcode{Certify} and 2.3~s for
a rejection followed by a kernel-checked refutation. A batch \mcode{Certify[\{\ldots\}]}
shares one Lean run, and verdicts are cached for the session by the theorem text, so
certifying the same thing twice costs nothing. \texttt{import Mathlib} adds 10--40~s,
which is why Tier~1 is opt-in. The exact engine, which needs no Lean, answers every
example in this chapter in milliseconds.}

% ---------------------------------------------------------------------------
\section{Using Mathilda from inside Lean}
\label{sec:cert-from-lean}

The same division of labour works in the other direction. The repository's
\texttt{lean/} directory is a small Lake package\index{Lean 4!tactics}, with only core
Lean and no Mathlib, whose tactics ask Mathilda for a certificate while a Lean user
writes a proof, and then check that certificate in Lean.

\begin{table}[ht]
\centering\small
\begin{tabular}{@{}p{0.27\linewidth}p{0.18\linewidth}p{0.22\linewidth}p{0.25\linewidth}@{}}
\toprule
Tactic & Goal & Mathilda supplies & Lean checks with \\
\midrule
\texttt{mathilda\_linear\_combination} & \(a=b\) from \(h_i: l_i=r_i\) & cofactors \(c_i\) with \(a-b=\sum c_i(l_i-r_i)\) & \texttt{grind}, rewriting \\
\texttt{mathilda\_sos} & \(0\le p\), \(p\le q\) from \(0\le g_i\) & triples \((d,\ell,i)\), \(q-p=\sum d\,\ell^2 g_i\) & hints \(0\le\ell^2\), \texttt{grind} \\
\texttt{mathilda\_prime} & \(\exists c,\ \texttt{PrattCertifies}\ p\ c\) & a Pratt certificate & \texttt{decide +kernel} \\
\texttt{\#mathilda "expr"} & none & \mcode{LeanForm} of the value & nothing; it proves nothing \\
\bottomrule
\end{tabular}
\caption{The tactics of the \texttt{lean/} package.}
\label{tab:cert-tactics}
\end{table}

A Lean file that uses them looks like this. It is Lean source from the package's
README, not Mathilda output, and it was not checked by this
book's build.

\begin{ccode}[Lean 4 source, illustrative (lean/README.md)]
import Mathilda

example (x y : Rat) (h1 : x^2 + y^2 = 1) (h2 : x = y) : 2*y^2 = 1 := by
  mathilda_linear_combination

example (x y : Rat) : 2*x*y ≤ x^2 + y^2 := by
  mathilda_sos

example : ∃ c, Mathilda.PrattCertifies 274177 c := by
  mathilda_prime

#mathilda "Factor[x^10-1]"
\end{ccode}

Each tactic serializes the goal and its hypotheses (local variables become
\texttt{v1}, \texttt{v2}, \ldots), starts \texttt{Mathilda -file} on a short script, and
reads back an answer. The server side is \mcode{Lean\textasciigrave
Internal\textasciigrave Oracle}, and it needs no Lean at all, so it can be shown here.

\begin{codepairs}
\pair{certified-answers/oracle/1}{The first example's query: the goal \(2v_2^2=1\),
the hypotheses \(v_1^2+v_2^2=1\) and \(v_1=v_2\). The answer is the list of cofactors.}
\pair{certified-answers/oracle/2}{The second: \(2v_1v_2\le v_1^2+v_2^2\). One square,
\(1\cdot(v_1-v_2)^2\), times the trivial multiplier \(g_0=1\).}
\pair{certified-answers/oracle/3}{\(0\le v_1^3\) under \(0\le v_1\): the square
\(v_1^2\) times the first hypothesis.}
\pair{certified-answers/oracle/4}{No certificate exists for a composite; the tactic
fails with Mathilda's reason.}
\pair{certified-answers/oracle/5}{What \texttt{\#mathilda} prints: an information
message, never a proof.}
\end{codepairs}

\calloutnote{Under the hood}{The answer is never elaborated as Lean syntax. The
package's \texttt{Mathilda/Answer.lean} is a hand-written parser, a pure function on
strings, that accepts only natural-number literals, the variables \texttt{v<i>}, the
operators \texttt{+ - * / \^{}} with a literal exponent, lists and tuples, with bounded
sizes (an answer of at most \(10^6\) characters, numerals of at most 2000 digits).
Anything else (\texttt{by}, \texttt{fun}, \texttt{sorry}, an unknown identifier) is
refused. The tactic builds terms from the parsed data with fixed quotations, at the
goal's own type, and the kernel checks the resulting proof. The package's test suite
runs the tactics against a fake Mathilda that lies, returning wrong cofactors,
\texttt{sorry}, and Lean code that tries to write a file, and checks that each one makes
the tactic fail without proving anything.}

% ---------------------------------------------------------------------------
\section{Limits}
\label{sec:cert-limits}

A certificate system earns trust by being precise about what it does not do.

\begin{itemize}
\item \textbf{Scope.} Tier~0 covers polynomial and rational identities, factorizations,
  antiderivatives built from rational functions, quadratic surds and one level of
  \B{Log}, \B{ArcTan} or \mcode{ArcTanh}, Bezout identities, ideal membership, rational
  roots, rational matrices (\B{Det} up to \(6\times6\)), sums of squares, polynomial
  optimization, primality, and closed-form sums of polynomial and geometric terms.
  Transcendental identities (\(\sin^2x+\cos^2x=1\)), \B{RootSum} and complex
  antiderivatives, radical roots of \B{Solve}, and hypergeometric sums such as
  \(\sum_k\binom nk^2\) have no certificate yet.
\item \textbf{Trusted steps.} Factor irreducibility, Risch non-elementarity,
  \B{NullSpace} completeness and Lucas's theorem are not proved. Each is named in
  \mcode{"TrustedSteps"} when it applies.
\item \textbf{Mathlib.} Analytic statements need Tier~1, a Lake project with Mathlib
  built for the same toolchain, and pay \(10\)--\(40\)~s for \texttt{import Mathlib}.
\item \textbf{The claim as written.} A claim that contains a computation other than
  the value-preserving rational ones, or a symbol with any kind of value, is refused.
  Certify the computation itself, or \mcode{Certify[Evaluate[claim]]} if the
  value-substituted claim is what you mean.
\item \textbf{\mcode{LeanVerify}'s surface.} It runs no source with user definitions,
  strings, comments or \texttt{\#} commands. A development of your own has to be
  checked by Lean and Lake directly.
\item \textbf{Exact is not kernel.} Without Lean, the best verdict available is
  \mcode{Missing["Exact", True]}, a re-check by the system that made the
  certificate.
\item \textbf{Time.} \texttt{system()} blocks, and \B{TimeConstrained} measures
  Mathilda's own CPU time, so it does not interrupt a Lean run; the bound is Lean's
  deterministic heartbeat count.
\end{itemize}

The next chapter returns to computation, but the habit this one teaches is worth
keeping. A CAS is most useful when you can tell which of its answers are proved,
which are checked, and which are only very probably right, and Mathilda now labels
each one.
